\section{The Three Interpretations Compared}
\label{sec:comparison}

All three interpretations take the householder as God, the denarius as eternal
life, and the first reading's contrast of thoughts as a statement about mercy.
All three find the same Christ in the parable: the one who went out to the
nations, the one who calls each person at every age, and the one whose
goodness is the reason for the wage. They differ in the question they put to
the parable, and therefore in the verses on which they rest and in what they
ask of the other texts.

\begin{comparisontable}{Question}{One wage across the ages}{The day as a human life}{Because I am good}
What is the day? & The history of the world, from the first just man to the
end. & One human life, from childhood to extreme old age. & Not decisive; the
evening's answer is.\\
Who are the last? & The Gentiles, and with them all who come after Christ
(Gregory); all Christians (Augustine); the people of the nations (Jerome's
report; Aquinas). & Those who turn to God late in life, and the thief on the
cross. & Anyone, as one who receives what is not owed.\\
Where does the weight fall? & On the hiring through the day, vv.~1--7. & On
the question at the eleventh hour, ``Why stand you here all the day idle?'',
v.~6. & On the householder's reply, vv.~13--15, from which the \work{Ordo}
takes its title.\\
What is the murmuring? & The long wait of the just before Christ (Gregory);
every earlier calling's envy of the Gentiles (Jerome). & A device of the
story, not to be pressed (Chrysostom). & The evil eye: grief at goodness (Aquinas); the
wish that another receive nothing (Jerome).\\
How is Isaiah 55 heard? & God's pardon reaches the nations that did not know
him (Jerome). & The time is short: while in the body (Jerome), before
adversity or death (Aquinas). & Mercy against requital (Aquinas); the
Lord most merciful and great in forgiving (Jerome).\\
What does Paul show? & The apostle of the nations at work for a church of the
last hour. & A life measured by its fruit (Chrysostom; Aquinas). & ``To live
is Christ'': he stays for Christ's honour and his people's progress (Aquinas;
Chrysostom).\\
What is asked at the Alleluia? & A place with Lydia, who signifies the Church
that received the preaching (Bede). & The hearing in which conversion begins.
& The opening Augustine counts God's gift; for Chrysostom God opens and Lydia
attends.\\
Principal witnesses & Gregory, Augustine, Aquinas; Jerome as reporter. &
Chrysostom, Jerome, Augustine, Gregory, Aquinas. & Augustine, Gregory, Jerome,
Aquinas, Hilary.\\
\end{comparisontable}

The first two interpretations are complementary. Augustine, Gregory and
Aquinas hold both and say so, and Jerome sets them side by side, preferring
the second. They answer different questions, how wide God's call has been and
how urgent it is for me, and neither excludes the other; Gregory passes from
one to the other in a single breath when he reminds his hearers that they, who
stand at the world's eleventh hour, were called at the first hour of their own
lives. Chrysostom stands apart. He holds the second interpretation alone and
gives the hours no reference to the ages of the world. Where Gregory and
Aquinas take the excuse ``No man hath hired us'' as the Gentiles' true plea,
he holds that God called all from the first. Yet he too hears in the saying
about first and last a dark hint at the Jews.

The third interpretation differs from the other two in kind. They read the
parable from its hours and it reads from the ending, and it is the only one of
the three for which the murmuring is the point and not a problem. Its
furthest step, that the first opening of the heart is God's gift, rests on
Augustine's argument. He and Chrysostom both speak of Acts 16:14 and read it
differently. Chrysostom says that God opens the hearts that are willing and
gives the opening to God and the attending to Lydia, without saying where the
willingness comes from; Augustine counts the very beginning of faith a gift
and says that Lydia was called so that she might believe. The difference
touches that step and does not reach the interpretation's centre, since
Chrysostom himself holds that in the parable ``the whole is of His love to
man''.

A second difference of exegesis crosses all three, and it concerns the
parable's last line. For Augustine ``the last shall be first'' states the
parable's equality; for Chrysostom it is a separate prophecy of reversal; for
Jerome it is the key by which the whole parable is read; Aquinas records the
first two. The Lectionary ends the Gospel on that line and leaves it open.

The Missal's prayers are patient of all three hearings because they say
something prior to each. The Lord is his people's salvation and hears them
whenever they cry; the law is love of God and neighbour, and its keeping ends
in life; what faith professes the sacrament gives; and redemption is to be
possessed at the altar and in conduct. A congregation that has been in the
vineyard since morning, a penitent who has come at the last hour, and a
labourer who has begun to count what he is owed all hear the same prayers
and receive the same Communion, the pledge of the one wage placed already in
every hand.
